\section{Integrals and Euler's Method}

\begin{enumerate}
  \item Consider the initial value problem
    \[
      y'(x)=f(x,y(x)), \qquad y(a)=y_0.
    \]
    Write down the forward Euler update on a partition
    \[
      a=x_0<x_1<\cdots<x_n=b
    \]
    and explain why one step has the form
    \[
      y_{i+1}-y_i \approx f(x_i,y_i)\,\Delta x_i.
    \]
  \item Why is a single Euler update naturally interpreted as an accumulation of a local rate over a small interval?
  \item Explain why it is reasonable to call
    \[
      y_{i+1}=y_i+f(x_i,y_i)\,\Delta x_i
    \]
    an integration step rather than merely an algebraic update.
  \item Sum the Euler increments over many steps and derive
    \[
      y_n
      =
      y_0+\sum_{i=0}^{n-1} f(x_i,y_i)\,\Delta x_i.
    \]
    What is the conceptual meaning of this formula, and how does the initial value~\(y_0\) enter as the starting accumulated amount?
  \item Why does the exact solution of
    \[
      y'(x)=f(x,y(x)), \qquad y(a)=y_0,
    \]
    satisfy the integral equation
    \[
      y(x)=y_0+\int_a^x f(t,y(t))\,dt?
    \]
    How does Euler's method mimic this identity?
  \item Why is the solution curve of an ODE often called an integral curve?
    What is being integrated, and in what sense is the curve determined by that integration?
  \item For a general first-order ODE
    \[
      y'=f(x,y),
    \]
    what information does the slope field encode?
    Why does prescribing an initial value
    \[
      y(a)=y_0
    \]
    select one particular curve from that field?
  \item In the autonomous case
    \[
      y'=F(y),
    \]
    how does the interpretation of an integral curve differ from the case
    \[
      y'=f(x,y)?
    \]
    What remains the same?
  \item Show that an ordinary definite integral is the special case of the ODE
    \[
      y'(x)=g(x), \qquad y(a)=y_0,
    \]
    where the right-hand side does not depend on~\(y\).
    What is the exact solution?
  \item Let~\(G\) be one antiderivative of~\(g\), so that
    \[
      G'(x)=g(x).
    \]
    Why are all solutions of
    \[
      y'(x)=g(x)
    \]
    of the form
    \[
      y(x)=G(x)+C?
    \]
    How does the initial condition
    \[
      y(a)=y_0
    \]
    determine the value of~\(C\)?
  \item Compare the slope field of a general ODE
    \[
      y'=f(x,y)
    \]
    with the slope field of
    \[
      y'=g(x).
    \]
    Why is the latter vertically uniform, meaning that at a fixed~\(x\) the slope is the same at every height~\(y\)?
  \item Why does the vertical uniformity of the slope field for
    \[
      y'=g(x)
    \]
    imply that all integral curves are vertical translates of one another?
    How does this connect to the additive constant~\(C\) in
    \[
      y=G(x)+C?
    \]
  \item Apply Euler's method to
    \[
      y'(x)=g(x), \qquad y(a)=y_0.
    \]
    Why does the scheme become
    \[
      y_{i+1}=y_i+g(x_i)\,\Delta x_i?
    \]
  \item Iterate the previous update and show that
    \[
      y_n
      =
      y_0+\sum_{i=0}^{n-1} g(x_i)\,\Delta x_i.
    \]
    Why is this now literally a Riemann sum plus the initial value, rather than only an Euler approximation?
  \item What special feature makes the discrete accumulation for
    \[
      y'=g(x)
    \]
    telescope exactly into
    \[
      y_n-y_0=\sum_{i=0}^{n-1} (y_{i+1}-y_i)?
    \]
    Why does the same telescoping identity not by itself turn a general Euler method for
    \[
      y'=f(x,y)
    \]
    into a classical Riemann sum for a known function of~\(x\) alone?
  \item In the case
    \[
      y'(x)=g(x),
    \]
    under what conditions on the partition and sample points does
    \[
      \sum_{i=0}^{n-1} g(x_i)\,\Delta x_i
    \]
    converge to
    \[
      \int_a^b g(x)\,dx?
    \]
    State the connection with the definition of Riemann integrability.
  \item If
    \[
      y'(x)=f(x,y)
    \]
    does depend on~\(y\), what breaks in the attempt to interpret Euler's method as an ordinary Riemann sum in the variable~\(x\) alone?
    Which additional unknown quantity is being approximated at every step?
  \item Compare these two viewpoints:
    \[
      \int_a^b g(x)\,dx
      \quad \text{and} \quad
      y(b)-y(a) \text{ for } y'=f(x,y).
    \]
    In what sense are both global accumulations of local data, and in what sense is the second problem structurally more implicit?
  \item If two solutions of
    \[
      y'=g(x)
    \]
    have different initial values, why does their difference remain constant in~\(x\)?
    Why is the analogous statement usually false for a general ODE
    \[
      y'=f(x,y)?
    \]
\end{enumerate}
